\chapter{模型扩展与工程实现建议}

\section{鲁棒估计扩展}
当观测包含粗差时，传统二范数目标函数对离群点敏感。可将目标函数改写为
\[
\min_{\bm p}\sum_i \rho\!\left(\frac{r_i(\bm p)}{\sigma_i}\right),
\]
其中 $\rho(\cdot)$ 取 Huber、Cauchy 或 Tukey 函数。此时可通过迭代重加权最小二乘（IRLS）求解，核心思想是根据残差大小动态更新权重，抑制异常值影响。

\section{动态场景扩展}
若目标随时间运动，静态位置估计可扩展为状态估计问题。令状态为
\[
\bm x_t=[x_t,y_t,z_t,\dot x_t,\dot y_t,\dot z_t]^{\T},
\]
建立状态转移
\[
\bm x_{t+1}=\bm F\bm x_t+\bm w_t
\]
和观测方程
\[
\bm y_t=\bm h(\bm x_t)+\bm v_t.
\]
可采用扩展卡尔曼滤波（EKF）或无迹卡尔曼滤波（UKF）实现实时跟踪。

\section{多目标扩展}
多目标场景需要解决“观测--目标”关联问题。常见方案包括：
\begin{itemize}
  \item 最近邻关联（NN）；
  \item 联合概率数据关联（JPDA）；
  \item 多假设跟踪（MHT）。
\end{itemize}
完成关联后，对每个目标独立应用本文定位求解器，或在联合优化框架中共享先验约束。

\section{约束优化扩展}
若目标存在先验约束（如高度区间、飞行走廊），可引入约束优化模型：
\[
\min_{\bm p} F(\bm p),\quad
\text{s.t.}\ g_\ell(\bm p)\le 0,\ h_s(\bm p)=0.
\]
可采用罚函数法、增广拉格朗日法或序列二次规划（SQP）。在小维问题中，简单边界投影也能显著改善异常估计。

\section{软件工程实现建议}
结合课程代码，给出以下重构建议：
\begin{enumerate}
  \item \textbf{模块分层}：观测模型层、导数层、求解器层、评估层独立；
  \item \textbf{统一数据结构}：用结构体封装站位、观测值、方差；
  \item \textbf{日志与可追踪}：每次迭代输出残差范数与步长；
  \item \textbf{异常处理}：矩阵奇异、角度奇异、数据缺失均返回明确状态码；
  \item \textbf{测试体系}：单元测试验证导数正确性，集成测试验证收敛结果。
\end{enumerate}

\section{参数选择经验}
\begin{itemize}
  \item 最大迭代次数：建议 $50\sim200$，并配合多重停机准则；
  \item 初值策略：优先几何初值，次选历史轨迹外推；
  \item 权重设置：尽量使用经过标定的标准差；
  \item 阈值设置：角度周期修正阈值建议与传感器定义保持一致。
\end{itemize}

\section{可视化建议}
为了提升结果可解释性，建议输出以下图形：
\begin{enumerate}
  \item 迭代轨迹图（$x$-$y$ 平面与 $z$ 方向）；
  \item 目标函数值与迭代次数关系；
  \item 三轴误差箱线图；
  \item 置信椭球三维可视化；
  \item 站位几何与目标相对分布图。
\end{enumerate}

\section{本章小结}
本章从鲁棒、动态、多目标和工程化四个维度给出扩展方向。相关建议可直接服务于后续课程深化项目或实际系统原型开发。
